\section{LRU Cache}
\label{sec:sq-lru-cache}

\problemstatement{Design a cache with a fixed capacity supporting
$\textit{get}(\textit{key})$ (return the value if present, else $-1$, and
mark the key as recently used) and $\textit{put}(\textit{key},
\textit{value})$ (insert or update a key, marking it recently used; if the
cache is full and a new key must be inserted, evict the
\textbf{least recently used} key first). Both operations must run in
$O(1)$.}

\begin{patternbox}
\emph{``Most/least recently used, with fast lookup''} is the keyword
combination from Section~\ref{sec:sq-intro}: $O(1)$ lookup by key demands
a hash map, while $O(1)$ ``move this to the most-recent end'' and
``evict the least-recent end'' demands a structure with fast removal and
insertion at \emph{both} ends and from the \emph{middle} -- which a plain
array or singly linked list cannot offer, but a \textbf{doubly linked
list} can.
\end{patternbox}

\subsection*{Understanding the problem carefully}

A hash map alone gives $O(1)$ lookup but no notion of ``recency
order.'' A queue alone gives recency order but $O(n)$ lookup by key, and
moving an arbitrary key to the front of recency would require knowing its
position. Combining both is the standard fix: a \textbf{doubly linked
list} maintains recency order (most recently used at one end, least
recently used at the other), and a \textbf{hash map from key to that
key's node} gives $O(1)$ lookup -- crucially, because the list is
\emph{doubly} linked, removing a node found via the hash map (to move it
or evict it) takes $O(1)$, since we have direct pointers to its neighbors
without needing to traverse from either end.

\begin{ideabox}[Doubly Linked List + Hash Map]
Maintain a doubly linked list with two sentinel nodes, \textit{head} and
\textit{tail}, where the side nearest \textit{head} represents
\emph{most} recently used and the side nearest \textit{tail} represents
\emph{least} recently used. A hash map stores $\textit{key} \to
\textit{node pointer}$. On $\textit{get}(\textit{key})$: if present,
unlink the node from its current position and re-insert it right after
\textit{head} (marking it most recent), then return its value; if absent,
return $-1$. On $\textit{put}(\textit{key}, \textit{value})$: if the key
exists, update its value and move its node to just after \textit{head}
(same as a successful get). If the key is new and the cache is at
capacity, remove the node just before \textit{tail} (the least recently
used) along with its hash map entry; then insert a new node for this key
just after \textit{head}.
\end{ideabox}

\subsection*{Why doubly linked, not singly linked, is essential}

The operation that makes this problem hard is not insertion or removal at
the \emph{ends} (a singly linked list with head and tail pointers could
handle that, as in Section~\ref{sec:sq-queue-linked-list}) -- it is
removal from an \textbf{arbitrary middle position}, which happens every
time an already-cached key is accessed again and must be promoted to
most-recently-used. Removing a node from the middle of a singly linked
list requires knowing its \emph{predecessor}, which means traversing from
the head -- an $O(n)$ operation. A doubly linked list stores each node's
predecessor pointer directly, so the hash map's stored node pointer alone
is enough to unlink it in $O(1)$, with no traversal needed at all.

\subsection*{Algorithm}

\begin{enumerate}
  \item Maintain sentinel \textit{head} and \textit{tail} nodes, a hash
        map $\textit{key} \to \textit{node}$, and a capacity limit.
  \item \textsc{Get}($\textit{key}$): if absent, return $-1$. Else:
        unlink the node; re-insert it right after \textit{head}; return
        its value.
  \item \textsc{Put}($\textit{key}, \textit{value}$): if present, update
        the node's value and move it to right after \textit{head}
        (as in \textsc{Get}). If absent and at capacity: remove the
        node just before \textit{tail} (and its map entry). Create a new
        node, insert it right after \textit{head}, and add it to the map.
\end{enumerate}

\subsection*{Dry run}

\dryrun{Capacity $2$. \textsc{Put}(1,1), \textsc{Put}(2,2),
\textsc{Get}(1), \textsc{Put}(3,3) (evicts LRU), \textsc{Get}(2).}

\begin{itemize}
  \item Put(1,1): list (MRU$\to$LRU) $=[1]$.
  \item Put(2,2): list $=[2,1]$.
  \item Get(1): found; move to front. List $=[1,2]$. Returns $1$.
  \item Put(3,3): cache full ($2$ keys); evict LRU end ($2$). List
        $=[3,1]$.
  \item Get(2): key $2$ was evicted; returns $-1$.
\end{itemize}

\subsection*{C++ implementation}

\begin{lstlisting}
#include <bits/stdc++.h>
using namespace std;

class LRUCache {
    struct Node {
        int key, value;
        Node* prev = nullptr;
        Node* next = nullptr;
        Node(int k, int v) : key(k), value(v) {}
    };

    int capacity;
    Node* head; // sentinel: head->next is most recently used
    Node* tail; // sentinel: tail->prev is least recently used
    unordered_map<int, Node*> map;

    void remove(Node* node) {
        node->prev->next = node->next;
        node->next->prev = node->prev;
    }

    void insertAfterHead(Node* node) {
        node->next = head->next;
        node->prev = head;
        head->next->prev = node;
        head->next = node;
    }

public:
    LRUCache(int cap) : capacity(cap) {
        head = new Node(-1, -1);
        tail = new Node(-1, -1);
        head->next = tail;
        tail->prev = head;
    }

    int get(int key) {
        if (!map.count(key)) return -1;
        Node* node = map[key];
        remove(node);
        insertAfterHead(node);
        return node->value;
    }

    void put(int key, int value) {
        if (map.count(key)) {
            Node* node = map[key];
            node->value = value;
            remove(node);
            insertAfterHead(node);
            return;
        }
        if ((int)map.size() == capacity) {
            Node* lru = tail->prev;
            remove(lru);
            map.erase(lru->key);
            delete lru;
        }
        Node* node = new Node(key, value);
        insertAfterHead(node);
        map[key] = node;
    }
};

int main() {
    LRUCache cache(2);
    cache.put(1, 1);
    cache.put(2, 2);
    cout << cache.get(1) << endl; // 1
    cache.put(3, 3); // evicts key 2
    cout << cache.get(2) << endl; // -1
    return 0;
}
\end{lstlisting}

\begin{complexitybox}
\textbf{Time Complexity.} $O(1)$ for both $\textit{get}$ and
$\textit{put}$ -- hash map lookup, plus a constant number of pointer
updates for the doubly linked list operations.
\textbf{Space Complexity.} $O(\textit{capacity})$ for the nodes and hash
map.
\end{complexitybox}

\begin{takeawaybox}
``Fast lookup by key'' plus ``fast reordering by recency, including
removal from the middle'' is the signature combination that demands a
hash map paired with a \emph{doubly} linked list, not a singly linked one
-- the doubly linked structure is what converts an $O(n)$ traversal-to-find-predecessor
into an $O(1)$ direct unlink. The same pairing, with a different
eviction rule, solves LFU Cache next.
\end{takeawaybox}
